\documentclass[aspectratio=169]{beamer}
\usetheme{Madrid}
\usecolortheme{default}
\usepackage[utf8]{inputenc}
\usepackage{amsmath, amssymb}
\usepackage{graphicx}
\usepackage{tcolorbox}

\definecolor{UCBblue}{RGB}{0, 51, 153}
\setbeamercolor{palette primary}{bg=UCBblue,fg=white}
\setbeamercolor{palette secondary}{bg=UCBblue,fg=white}
\setbeamercolor{palette tertiary}{bg=UCBblue,fg=white}
\setbeamercolor{titlelike}{bg=UCBblue,fg=white}
\setbeamercolor{item}{fg=UCBblue}

\title[Nivelación - Sesión 2]{Taller de Nivelación Intensiva:\\ Álgebra Lineal Aplicada a la Teoría de Redes}
\subtitle{Módulo 3: El Límite de la CC y el Análisis Fasorial}
\author[M.Sc. Ing. B. Quiroga Turdera]{M.Sc. Ing. Bernardo Quiroga Turdera}
\institute[UCB Tarija]{Carrera de Ingeniería Mecatrónica\\ Universidad Católica Boliviana "San Pablo"}
\date{29 de Julio de 2026}

\begin{document}

\begin{frame}
    \titlepage
\end{frame}

\begin{frame}{Agenda de la Sesión}
    \tableofcontents
\end{frame}

\section{El Límite de la Corriente Continua (CC)}
\begin{frame}{¿Por qué la Corriente Continua no es suficiente?}
    Hasta ahora, hemos modelado sistemas en \textbf{Corriente Continua (CC)}, donde la energía fluye de manera constante e invariable en el tiempo.
    
    \vspace{0.3cm}
    \begin{alertblock}{El Problema Industrial}
    La inmensa mayoría de los actuadores industriales (motores de inducción), la red eléctrica global y las señales de los sensores operan en \textbf{Corriente Alterna (CA)}. En CA, el voltaje y la corriente son ondas senoidales que cambian de magnitud y polaridad cíclicamente.
    \end{alertblock}
    
    \vspace{0.2cm}
    \textit{Consecuencia analítica:} Una matriz de resistencias puras ($\mathbf{R}$) ya no puede capturar la dinámica de componentes que almacenan energía (bobinas y capacitores).
\end{frame}

\section{El Plano Complejo y la Impedancia}
\begin{frame}{La Evolución de la Resistencia: La Impedancia ($\mathbf{Z}$)}
    En CA, la oposición al flujo de electrones tiene dos dimensiones físicas que se modelan en el \textbf{Plano Complejo ($s$)}:
    
    \vspace{0.3cm}
    \begin{itemize}
        \item \textbf{Parte Real ($R$):} La resistencia clásica que disipa energía en forma de calor.
        \item \textbf{Parte Imaginaria ($X$):} La \textit{reactancia}. Representa la energía que se almacena y se devuelve al circuito (campos magnéticos en bobinas, campos eléctricos en capacitores).
    \end{itemize}
    
    \vspace{0.3cm}
    \begin{tcolorbox}[colback=blue!5!white,colframe=UCBblue,title=Definición Matemática de la Impedancia]
    $$ \mathbf{Z} = R + jX $$
    Donde $j$ es el operador imaginario ($\sqrt{-1}$). En ingeniería eléctrica usamos $j$ en lugar de $i$ para no confundirlo con la variable de corriente.
    \end{tcolorbox}
\end{frame}

\section{El Concepto de Fasor}
\begin{frame}{De la Onda en el Tiempo al Vector en el Espacio}
    Manejar senoides puras ($v(t) = V_m \cos(\omega t + \theta)$) en sistemas de ecuaciones matriciales genera cálculos trigonométricos intratables.
    
    \vspace{0.3cm}
    \textbf{La Solución de Ingeniería (Transformada Fasorial):}
    Congelamos el tiempo y representamos la onda como un vector estático (Fasor) que tiene dos propiedades fundamentales:
    \begin{enumerate}
        \item \textbf{Magnitud ($|\mathbf{Z}|$):} La longitud del vector (cuánta oposición total existe).
        \item \textbf{Fase ($\theta$):} El ángulo del vector (el desfase temporal entre el voltaje y la corriente).
    \end{enumerate}
\end{frame}

\section{Transición al Gemelo Digital}
\begin{frame}{Sintaxis en Python para el Dominio Fasorial}
    \textbf{Fase CDIO: Implementar}
    
    \vspace{0.3cm}
    Python está diseñado nativamente para el cómputo científico y soporta números complejos sin librerías externas adicionales para su declaración:
    \begin{itemize}
        \item \textbf{Operador Imaginario:} En Python, se utiliza el sufijo \texttt{j} adyacente al número (Ej. \texttt{1j}, \texttt{5+2j}).
        \item \textbf{Magnitud (Forma Polar):} Extraída computacionalmente con \texttt{np.abs(Z)}.
        \item \textbf{Fase o Ángulo:} Extraída con \texttt{np.angle(Z, deg=True)}.
    \end{itemize}
    
    \vspace{0.5cm}
    \begin{center}
        \textit{Pasemos a la Guía Analítica para comprender la conversión trigonométrica antes de programar.}
    \end{center}
\end{frame}

\end{document}
